% GENERATED FILE. Edit canonical lessons, then run npm run generate:books.
% canonical-source-hash: fnv1a64:fe645ebc0c51afcf
% canonical-lessons: FR-C373-disponible, FR-C373-engager, FR-C373-engagement, FR-C373-flexible, FR-C373-envisager, FR-R373-revision-366-373

\chapter{Si j'ai le temps}
\label{ch:fr-si-j-ai-le-temps}
\hlchaptermodality{373}

\begin{chapteropening}
\textbf{By the end of this chapter:} \emph{I can say whether I am free, commit to something, say I am flexible, and say what I am considering.}

Complete the last lesson of chapter 373: I can say whether I am free, commit to something, say I am flexible, and say what I am considering.
\end{chapteropening}

\section[disponible]{disponible — available, free}

\label{lesson:FR-C373-disponible}



\begin{quote}
Before the new one: say the French for the probability, then the French for even if it means.
\end{quote}

\subsection*{You'll want to know: disponible}
\textbf{disponible} — \textquotedblleft{}available, free\textquotedblright{}.

\begin{quote}
Je suis disponible mardi après-midi.
\end{quote}

\begin{quote}
\emph{I am free on Tuesday afternoon.}
\end{quote}

Of people and things: \emph{ce modèle n'est plus disponible}.

\begin{cousinweb}
\emph{disponible} is a learned borrowing from Medieval Latin \emph{dispōnibilis} (EtymDB), from Latin \emph{dispōnere} \textquotedblleft{}to set in order, to arrange\textquotedblright{}, built on \emph{pōnere} \textquotedblleft{}to put\textquotedblright{} (Lewis \& Short).
\end{cousinweb}

\begin{grammarlens}[title={What you've built}]
One more word for telling, explaining and describing at B1.
\end{grammarlens}

\subsection*{Your turn}
\begin{itemize}
\raggedright
  \item \emph{Say it:} \emph{disponible}
  \item \emph{Say it:} \emph{disponible}, once more
  \item \emph{Say it:} a sentence of your own with \emph{disponible}
  \item \emph{Recall:} say the French for the probability, then the French for even if it means, then say \emph{disponible} again
\end{itemize}

\begin{tcolorbox}[breakable,colback=teal!4,colframe=teal!35!black,title={Before you move on}]
What does disponible mean? (\textquotedblleft{}Available, free\textquotedblright{}.) Say it once more.
\end{tcolorbox}
\section[engager]{engager — to commit (mostly reflexive, s'engager); to hire}

\label{lesson:FR-C373-engager}



\begin{quote}
Before the new one: say the French for even if it means, then the French for available.
\end{quote}

\subsection*{You'll want to know: engager}
\textbf{engager} — \textquotedblleft{}to commit (mostly reflexive, s'engager); to hire\textquotedblright{}.

\begin{quote}
Il s'est engagé à finir le travail avant vendredi.
\end{quote}

\begin{quote}
\emph{He committed to finishing the work before Friday.}
\end{quote}

\emph{S'engager à} + infinitive — to commit to doing. \emph{Engager quelqu'un} — to hire someone. \emph{Nous engageons}.

\begin{cousinweb}
\emph{engager} is built from \emph{en-} and \emph{gage} \textquotedblleft{}a pledge\textquotedblright{}. \emph{gage} is usually traced to a Frankish word for a pledge (EtymDB lists Germanic relatives such as Gothic \emph{wadi}).
\end{cousinweb}

\begin{grammarlens}[title={What you've built}]
One more word for telling, explaining and describing at B1.
\end{grammarlens}

\subsection*{Your turn}
\begin{itemize}
\raggedright
  \item \emph{Say it:} \emph{engager}
  \item \emph{Say it:} \emph{engager}, once more
  \item \emph{Say it:} a sentence of your own with \emph{engager}
  \item \emph{Recall:} say the French for even if it means, then the French for available, then say \emph{engager} again
\end{itemize}

\begin{tcolorbox}[breakable,colback=teal!4,colframe=teal!35!black,title={Before you move on}]
What does engager mean? (\textquotedblleft{}To commit (mostly reflexive, s'engager); to hire\textquotedblright{}.) Say it once more.
\end{tcolorbox}
\section[l'engagement]{l'engagement — the commitment}

\label{lesson:FR-C373-engagement}



\begin{quote}
Before the new one: say the French for available, then the French for to commit.
\end{quote}

\subsection*{You'll want to know: l'engagement}
\textbf{l'engagement} — \textquotedblleft{}the commitment\textquotedblright{}.

\begin{quote}
C'est un engagement important : je ne peux pas dire non.
\end{quote}

\begin{quote}
\emph{It is an important commitment: I cannot say no.}
\end{quote}

\emph{Prendre un engagement} — to make a commitment. Masculine.

\begin{cousinweb}
\emph{engagement} is built from \emph{engager} and \emph{-ment} (EtymDB). \emph{engager} rests on \emph{gage} \textquotedblleft{}a pledge\textquotedblright{}, usually traced to a Frankish word.
\end{cousinweb}

\begin{grammarlens}[title={What you've built}]
One more word for telling, explaining and describing at B1.
\end{grammarlens}

\subsection*{Your turn}
\begin{itemize}
\raggedright
  \item \emph{Say it:} \emph{l'engagement}
  \item \emph{Say it:} \emph{l'engagement}, once more
  \item \emph{Say it:} a sentence of your own with \emph{l'engagement}
  \item \emph{Recall:} say the French for available, then the French for to commit, then say \emph{l'engagement} again
\end{itemize}

\begin{tcolorbox}[breakable,colback=teal!4,colframe=teal!35!black,title={Before you move on}]
What does l'engagement mean? (\textquotedblleft{}The commitment\textquotedblright{}.) Say it once more.
\end{tcolorbox}
\section[flexible]{flexible — flexible}

\label{lesson:FR-C373-flexible}



\begin{quote}
Before the new one: say the French for to commit, then the French for the commitment.
\end{quote}

\subsection*{You'll want to know: flexible}
\textbf{flexible} — \textquotedblleft{}flexible\textquotedblright{}.

\begin{quote}
Mon patron est assez flexible sur les horaires.
\end{quote}

\begin{quote}
\emph{My boss is quite flexible about hours.}
\end{quote}

One form for both genders. Close to \emph{souple}.

\begin{cousinweb}
\emph{flexible} is a learned borrowing from Latin \emph{flexibilis} \textquotedblleft{}easily bent\textquotedblright{} (EtymDB; Lewis \& Short), from \emph{flectere} \textquotedblleft{}to bend\textquotedblright{}.
\end{cousinweb}

\begin{grammarlens}[title={What you've built}]
One more word for telling, explaining and describing at B1.
\end{grammarlens}

\subsection*{Your turn}
\begin{itemize}
\raggedright
  \item \emph{Say it:} \emph{flexible}
  \item \emph{Say it:} \emph{flexible}, once more
  \item \emph{Say it:} a sentence of your own with \emph{flexible}
  \item \emph{Recall:} say the French for to commit, then the French for the commitment, then say \emph{flexible} again
\end{itemize}

\begin{tcolorbox}[breakable,colback=teal!4,colframe=teal!35!black,title={Before you move on}]
What does flexible mean? (\textquotedblleft{}Flexible\textquotedblright{}.) Say it once more.
\end{tcolorbox}
\section[envisager]{envisager — to consider, to plan on}

\label{lesson:FR-C373-envisager}



\begin{quote}
Before the new one: say the French for the commitment, then the French for flexible.
\end{quote}

\subsection*{You'll want to know: envisager}
\textbf{envisager} — \textquotedblleft{}to consider, to plan on\textquotedblright{}.

\begin{quote}
Nous envisageons de déménager l'année prochaine.
\end{quote}

\begin{quote}
\emph{We are considering moving next year.}
\end{quote}

\emph{Envisager de} + infinitive. \emph{Nous envisageons}.

\begin{cousinweb}
\emph{envisager} is built from \emph{en-} and \emph{visage} \textquotedblleft{}face\textquotedblright{}: to look something in the face. \emph{visage} rests on Latin \emph{vīsus} \textquotedblleft{}sight\textquotedblright{}, from \emph{vidēre} \textquotedblleft{}to see\textquotedblright{} (Lewis \& Short).
\end{cousinweb}

\begin{grammarlens}[title={What you've built}]
One more word for telling, explaining and describing at B1.
\end{grammarlens}

\subsection*{Your turn}
\begin{itemize}
\raggedright
  \item \emph{Say it:} \emph{envisager}
  \item \emph{Say it:} \emph{envisager}, once more
  \item \emph{Say it:} a sentence of your own with \emph{envisager}
  \item \emph{Recall:} say the French for the commitment, then the French for flexible, then say \emph{envisager} again
\end{itemize}

\begin{tcolorbox}[breakable,colback=teal!4,colframe=teal!35!black,title={Before you move on}]
What does envisager mean? (\textquotedblleft{}To consider, to plan on\textquotedblright{}.) Say it once more.
\end{tcolorbox}
\section[(dialogue)]{Retour sur les chapitres 366–373}

\label{lesson:FR-R373-revision-366-373}



\begin{quote}
Say the French for the volunteer and the volunteering.
\end{quote}

\subsection*{Your turn}
Say these aloud:

\begin{itemize}
\raggedright
  \item the volunteer, the volunteering, the association, the solidarity, the community
  \item to suppose, the hypothesis, to presume, plausible, the supposition
  \item provided that, as long as, unless, in case, otherwise
  \item according to, possible, possibly, the possibility, for lack of
  \item the hitch, foreseeable, unpredictable, to adapt, supple
  \item to cancel, to postpone, temporary, final, to schedule
  \item the risk, risky, improbable, the probability, even if it means
  \item available, to commit, the commitment, flexible, to consider
\end{itemize}

\begin{tcolorbox}[breakable,colback=teal!4,colframe=teal!35!black,title={Before you move on}]
The volunteer? (\textbf{le bénévole}.) The volunteering? (\textbf{le bénévolat}.) The association? (\textbf{l'association}.) The solidarity? (\textbf{la solidarité}.) The community? (\textbf{la communauté}.) To suppose? (\textbf{supposer}.) The hypothesis? (\textbf{l'hypothèse}.) To presume? (\textbf{présumer}.) Plausible? (\textbf{plausible}.) The supposition? (\textbf{la supposition}.) Provided that? (\textbf{à condition que}.) As long as? (\textbf{pourvu que}.) Unless? (\textbf{à moins que}.) In case? (\textbf{au cas où}.) Otherwise? (\textbf{sinon}.) According to? (\textbf{selon}.) Possible? (\textbf{éventuel}.) Possibly? (\textbf{éventuellement}.) The possibility? (\textbf{l'éventualité}.) For lack of? (\textbf{faute de}.) The hitch? (\textbf{l'imprévu}.) Foreseeable? (\textbf{prévisible}.) Unpredictable? (\textbf{imprévisible}.) To adapt? (\textbf{adapter}.) Supple? (\textbf{souple}.) To cancel? (\textbf{annuler}.) To postpone? (\textbf{reporter}.) Temporary? (\textbf{provisoire}.) Final? (\textbf{définitif}.) To schedule? (\textbf{programmer}.) The risk? (\textbf{le risque}.) Risky? (\textbf{hasardeux}.) Improbable? (\textbf{improbable}.) The probability? (\textbf{la probabilité}.) Even if it means? (\textbf{quitte à}.) Available? (\textbf{disponible}.) To commit? (\textbf{engager}.) The commitment? (\textbf{l'engagement}.) Flexible? (\textbf{flexible}.) To consider? (\textbf{envisager}.)
\end{tcolorbox}
